\section{Settle the argument before drafting}
\label{sec:architecture}

The section list should follow the logic of the evidence. A fixed outline can
hide the paper's strongest result under routine setup or give equal space to
claims of very different importance. Decide the reader's prerequisites and the
evidence for each step. Section titles and paragraph topics follow from that
sequence.

\subsection{Open with the scientific bottleneck}

Begin with the scientific goal and the resource or assumption that blocks it.
Do not let the nearest acronym or prior method define the motivation unless a
reader must understand that method to see the problem. A method-specific
opening can make a broad question look like a narrow competition between
techniques.

Name the scarce and available resources. For example, a paper about robot
adaptation might study a setting with few expert
demonstrations but ample autonomous interaction. That framing identifies the
demonstration bottleneck before choosing imitation learning, in-context
learning, offline learning, or online reinforcement learning as the response.

\subsection{Choose the structure from the evidence}

The example papers use three recurring architectures.

\paragraph{Empirical method papers.}
\emph{DPPO}~\citep{ren2024diffusion} and
\emph{OGPO}~\citep{patil2026ogpo} begin with an operational bottleneck,
identify a tradeoff in current methods, and introduce a method that changes
that tradeoff. Headline comparisons establish the main result. Later
experiments test the proposed explanation. They also identify performance
limits and test deployment. This architecture works when performance is the
central claim but mechanism and reliability determine the result's
credibility.

\paragraph{Empirical mechanism papers.}
\emph{Much Ado About Noising}~\citep{pan2025muchado} organizes the paper
around competing explanations. Controlled tests first rule out explanations
that do not account for the observed performance. A minimal replacement then
isolates the components that do matter, and later experiments examine their
mechanism. This architecture fits a paper whose contribution is a corrected
understanding. The negative and positive findings must remain distinct:
rejecting one explanation does not prove another.

\paragraph{Theory-led papers.}
\emph{Action Chunking}~\citep{zhang2025actionchunking} starts from a failure
in continuous-control behavior cloning and studies two interventions in
different regimes. The formal results establish each intervention's sufficient
regime, while experiments test the predictions that matter outside the
model. This architecture fits a paper that moves from a failure or
impossibility to guarantees under stated conditions. The regime should appear
before the guarantee, both in prose and in section titles.

Treat these as shapes. A new paper may need a different order. All four
examples open with a concrete problem. Each then exposes a mismatch in current
thinking and resolves it as far as the evidence allows. Each section should
resolve a question created by the previous one.

\subsection{Present the decisive result first}

In an empirical paper, show the headline result before low-level ablations.
The reader needs to know that the method changes the outcome before learning
the relative importance of regularizers and schedules. Mechanism studies
should follow the performance claim unless mechanism is itself the main claim.

In a theory paper, present the motivating failure before the intervention.
State assumptions before the result that uses them, but avoid a long block of
notation that will not be needed for several pages. Definitions are
easier to absorb when they arrive close to the question they answer.

Related work can appear early or late. Place it early when the reader must
distinguish the method from close alternatives before understanding the
formulation. Place it after the main results when a literature survey would
delay the paper's question and answer. In either case, organize related work by
technical distinction. A chronological list of papers does not identify the
unresolved technical point.

\subsection{Draft the evidence before the introduction}

The order of the paper need not be the order of writing. A reliable drafting
sequence is:

\begin{enumerate}[leftmargin=*]
    \item settle the claim--evidence map;
    \item draft the central figures, tables, theorem statements, and captions;
    \item write the results that interpret that evidence;
    \item write the setup and method needed to understand those results;
    \item write the introduction once the contribution hierarchy is stable;
    \item write the abstract and title last.
\end{enumerate}

This sequence reduces a common source of churn. An early abstract often
commits the paper to claims unsupported by the final evidence. An early
introduction can also turn provisional organization into a constraint. The
abstract should explain the argument supported by the finished evidence, not
tour the paper's sections or describe the paper the authors first imagined.
